%! Statistics — Unit 1, Lesson 1.3 — Slides (Beamer)
\documentclass[aspectratio=169, 11pt]{beamer}
\usepackage{apstats-beamer}

% ── Deck-local layout ────────────────────────────────────────────────────────
% One left edge for the whole deck: the body text block and the header-bar
% title both start 0.9cm from the page edge, so nothing hangs out of line.
\setbeamersize{text margin left=0.9cm, text margin right=0.9cm}

% \slideheader{bar color}{Title} — same geometry as \navyheader (1.35cm bar,
% title on its vertical centre) but with the title inset matched to the text
% margin above. Used on every content slide so the header sits identically on
% each one, whatever the accent color.
\newcommand{\slideheader}[2]{%
  \begin{tikzpicture}[remember picture, overlay]
    \fill[#1] (current page.north west)
      rectangle ([yshift=-1.35cm]current page.north east);
    \node[anchor=west, text=white, font=\bfseries\large,
          inner xsep=0.9cm, inner ysep=0pt]
      at ([yshift=-0.68cm]current page.north west) {#2};
  \end{tikzpicture}%
  \vspace{1.15cm}\par}

%=============================================================================
\begin{document}

% ─────────────────────────────────────────────────────────────────────────────
% SLIDE 1 — LESSON TITLE
% ─────────────────────────────────────────────────────────────────────────────
{
\setbeamercolor{background canvas}{bg=navy}
\begin{frame}[plain]
  \begin{tikzpicture}[remember picture, overlay]
    \fill[navylight] (current page.north west)
      rectangle ([yshift=-1.35cm]current page.north east);
  \end{tikzpicture}
  \vspace{2.8cm}
  \begin{minipage}[t]{\textwidth}
    {\color{goldacc}\bfseries\small LESSON 1.3}\\[0.5cm]
    {\color{white}\bfseries\LARGE Representing a Categorical Variable\\[0.25cm]
      with Tables}\\[1.4cm]
    {\color{skymid}\small Statistics~$\cdot$~Unit 1: Exploring One-Variable Data}
  \end{minipage}
\end{frame}
}

% ─────────────────────────────────────────────────────────────────────────────
% SLIDE 2 — WARM-UP
% ─────────────────────────────────────────────────────────────────────────────
{
\setbeamercolor{background canvas}{bg=sky}
\begin{frame}[t, plain]
  \slideheader{navy}{Warm-Up}
  \begin{minipage}[t]{\textwidth}
    \small

    \begin{enumerate}
      \setlength{\itemsep}{0.7em}\setlength{\topsep}{0em}
      \item Look at the list of 20 students' favorite subjects on your warm-up sheet.
            Can you immediately tell which subject is most popular? Why or why not?
      \item Tally each subject and complete the frequency column.
            Does your total equal 20?
      \item From Lesson 1.2: what \textit{type} of variable is ``favorite subject''?
            Is it meaningful to compute the average?
      \item What question do you think today's lesson will answer?
    \end{enumerate}

    \vspace{0.6cm}
    {\color{navy}\itshape\small
      We'll revisit your frequency table at the end of the lesson to add
      one more column.}
  \end{minipage}
\end{frame}
}

% ─────────────────────────────────────────────────────────────────────────────
% SLIDE 3 — PRIMARY OBJECTIVE
% ─────────────────────────────────────────────────────────────────────────────
{
\setbeamercolor{background canvas}{bg=sky}
\begin{frame}[t, plain]
  \begin{tikzpicture}[remember picture, overlay]
    \fill[navy] (current page.north west)
      rectangle ([yshift=-0.25cm]current page.north east);
  \end{tikzpicture}
  \vspace{0.45cm}\par
  \begin{minipage}[t]{\textwidth}

    \sectionlabel{Primary Objective}

    \normalsize
    Students will construct and interpret \textbf{frequency} and
    \textbf{relative frequency} tables for a categorical variable.
    Students will explain why relative frequency --- not raw count --- is
    required to compare groups of different sizes (Big Idea: UNC).

    \vspace{0.45cm}

    \normalsize\textbf{\color{navy}By the end of this lesson, I will be able to:}

    \smallskip
    \begin{itemize}
      \setlength{\itemsep}{0.3em}\setlength{\topsep}{0.2em}
      \item Construct a frequency and relative frequency table from raw categorical data.
      \item Interpret relative frequencies in context using complete sentences.
      \item Explain why relative frequency is necessary for fair comparisons.
    \end{itemize}

  \end{minipage}
\end{frame}
}

% ─────────────────────────────────────────────────────────────────────────────
% SLIDE 4 — HOOK
% ─────────────────────────────────────────────────────────────────────────────
{
\setbeamercolor{background canvas}{bg=hookbg}
\begin{frame}[t, plain]
  \slideheader{goldacc}{Hook --- Which School Is More of a Walking School?}
  \begin{minipage}[t]{\textwidth}
    \small
    {\color{white}\bfseries
      Two schools survey their students about how they get to school.}

    \vspace{0.4cm}

    \setlength{\tabcolsep}{10pt}
    \renewcommand{\arraystretch}{1.35}
    \hspace*{\tabcolsep}\begin{tabular}{@{} >{\bfseries\color{white}}p{0.20\textwidth}
                        >{\centering\arraybackslash}p{0.20\textwidth}
                        >{\centering\arraybackslash}p{0.20\textwidth} @{}}
      \rowcolor{navy}
      {\color{skymid}Mode} &
      {\color{skymid}School A} &
      {\color{skymid}School B} \\
      \rowcolor{white!15!hookbg}
      {\color{white}Bus}  & {\color{white}120} & {\color{white}30} \\
      \rowcolor{white!8!hookbg}
      {\color{white}Car}  & {\color{white}80}  & {\color{white}50} \\
      \rowcolor{white!15!hookbg}
      {\color{white}Walk} & {\color{white}40}  & {\color{white}120} \\
      \rowcolor{navy}
      {\color{skymid}\textbf{Total}} &
      {\color{skymid}\textbf{240}} &
      {\color{skymid}\textbf{200}} \\
    \end{tabular}

    \vspace{0.45cm}
    {\color{white}\small
      Discuss with your partner: do the two Walk counts, 40 and 120, answer
      that question on their own?}
  \end{minipage}
\end{frame}
}

% ─────────────────────────────────────────────────────────────────────────────
% SLIDE 5 — HOOK REVEAL
% ─────────────────────────────────────────────────────────────────────────────
{
\setbeamercolor{background canvas}{bg=hookbg}
\begin{frame}[t, plain]
  \slideheader{goldacc}{Hook --- Counts Say ``How Many,'' Shares Say ``How Common''}
  \begin{minipage}[t]{\textwidth}
    \small

    \setlength{\tabcolsep}{8pt}
    \renewcommand{\arraystretch}{1.4}
    \hspace*{\tabcolsep}\begin{tabular}{@{} >{\bfseries\color{white}}p{0.14\textwidth}
                        >{\centering\arraybackslash}p{0.13\textwidth}
                        >{\centering\arraybackslash}p{0.17\textwidth}
                        >{\centering\arraybackslash}p{0.13\textwidth}
                        >{\centering\arraybackslash}p{0.17\textwidth} @{}}
      \rowcolor{navy}
      {\color{skymid}Mode} &
      {\color{skymid}Sch.\ A} & {\color{skymid}Rel.\ Freq.\ A} &
      {\color{skymid}Sch.\ B} & {\color{skymid}Rel.\ Freq.\ B} \\
      \rowcolor{white!15!hookbg}
      {\color{white}Bus}  & {\color{white}120} & {\color{white}50\%} &
        {\color{white}30}  & {\color{white}15\%} \\
      \rowcolor{white!8!hookbg}
      {\color{white}Car}  & {\color{white}80}  & {\color{white}33\%} &
        {\color{white}50}  & {\color{white}25\%} \\
      \rowcolor{navy}
      {\color{white}\textbf{Walk}} &
        {\color{white}40}  & {\color{goldacc}\textbf{17\%}} &
        {\color{white}120} & {\color{goldacc}\textbf{60\%}} \\
    \end{tabular}

    \vspace{0.45cm}
    {\color{white}\small
      The counts, 40 and 120, come out of different-sized schools ($n = 240$
      and $n = 200$), so they do not answer the question on their own.\\[6pt]
      The shares do: \textbf{60\%} of School B walks, against \textbf{17\%} of
      School A.\\[6pt]
      \textbf{Lesson:} divide by the total before you compare two groups.}
  \end{minipage}
\end{frame}
}

% ─────────────────────────────────────────────────────────────────────────────
% SLIDE 6 — VOCABULARY
% ─────────────────────────────────────────────────────────────────────────────
{
\setbeamercolor{background canvas}{bg=greenbg}
\setbeamercolor{itemize item}{fg=greenacc}
\begin{frame}[t, plain]
  \slideheader{navy}{Vocabulary}
  \begin{minipage}[t]{\textwidth}
    \small
    \begin{itemize}
      \setlength{\itemsep}{0.5em}\setlength{\topsep}{0em}
      \item {\color{navy}\bfseries Frequency} --- the count of individuals in a
            particular category.
      \item {\color{navy}\bfseries Relative frequency} --- frequency $\div$ total;
            expressed as a proportion (decimal) or a percent.
      \item {\color{navy}\bfseries Frequency table} --- a table listing each category
            of a categorical variable and its count.
      \item {\color{navy}\bfseries Relative frequency table} --- a table listing each
            category and its relative frequency.
      \item {\color{navy}\bfseries Distribution (categorical)} --- the set of possible
            values and their associated relative frequencies.
    \end{itemize}

    \vspace{0.55cm}
    {\color{navy}\small\textit{Key check: relative frequencies must sum to 1 (or 100\%).}}
  \end{minipage}
\end{frame}
}

% ─────────────────────────────────────────────────────────────────────────────
% SLIDE 7 — BUILDING A TABLE
% ─────────────────────────────────────────────────────────────────────────────
{
\setbeamercolor{background canvas}{bg=white}
\begin{frame}[t, plain]
  \slideheader{navy}{Building a Relative Frequency Table --- Four Steps}
  \begin{minipage}[t]{\textwidth}
    \small

    \begin{enumerate}
      \setlength{\itemsep}{0.55em}\setlength{\topsep}{0em}
      \item \textbf{List all categories} in the first column.
      \item \textbf{Tally} the raw data and record the \textbf{frequency} (count).
      \item \textbf{Divide} each count by $n$ (the total) to get the
            \textbf{relative frequency}.
      \item \textbf{Check:} all relative frequencies must sum to exactly 1.
    \end{enumerate}

    \vspace{0.4cm}
    \[
      \text{Relative frequency} = \frac{\text{frequency}}{n}
    \]

    \vspace{0.4cm}
    {\color{navy}\small
      Express as a decimal \textit{or} a percent --- both are correct.}
  \end{minipage}
\end{frame}
}

% ─────────────────────────────────────────────────────────────────────────────
% SLIDE 8 — WORKED EXAMPLE
% ─────────────────────────────────────────────────────────────────────────────
{
\setbeamercolor{background canvas}{bg=white}
\begin{frame}[t, plain]
  \slideheader{navy}{Worked Example --- Pet Ownership ($n = 40$)}
  \begin{minipage}[t]{\textwidth}
    \small

    \setlength{\tabcolsep}{10pt}
    \renewcommand{\arraystretch}{1.4}
    \hspace*{\tabcolsep}\begin{tabular}{@{} >{\bfseries}p{0.22\textwidth}
                        >{\centering\arraybackslash}p{0.16\textwidth}
                        >{\centering\arraybackslash}p{0.24\textwidth} @{}}
      \rowcolor{navy}
      {\color{white}Pet type} & {\color{white}Frequency} &
        {\color{white}Rel.\ frequency} \\
      Dog   & 12 & $12/40 = 0.30$ \\
      Cat   & 8  & $8/40  = 0.20$ \\
      Fish  & 6  & $6/40  = 0.15$ \\
      Bird  & 4  & $4/40  = 0.10$ \\
      None  & 10 & $10/40 = 0.25$ \\
      \rowcolor{sky} Total & 40 & 1.00 \\
    \end{tabular}

    \vspace{0.5cm}
    {\color{navy}\bfseries\small Model interpretation:}\\[4pt]
    {\small\itshape ``Of the 40 households surveyed, the most common pet type was
    dog, owned by 30\% of households. The least common was bird, owned by 10\%.''}
  \end{minipage}
\end{frame}
}

% ─────────────────────────────────────────────────────────────────────────────
% SLIDE 9 — WRITING THE DESCRIPTION
% ─────────────────────────────────────────────────────────────────────────────
{
\setbeamercolor{background canvas}{bg=sky}
\begin{frame}[t, plain]
  \slideheader{navy}{Describing a Distribution --- Saying It Precisely}
  \begin{minipage}[t]{\textwidth}
    \small

    \textbf{\color{navy}Every description must include:}
    \begin{itemize}
      \setlength{\itemsep}{0.35em}\setlength{\topsep}{0.2em}
      \item The \textbf{total $n$} (``of the 40 households surveyed\ldots'')
      \item The \textbf{variable} in context (``pet type'')
      \item The \textbf{most} and \textbf{least} common categories
      \item The \textbf{relative frequency} for at least the key categories
    \end{itemize}

    \vspace{0.45cm}
    {\color{navy}\bfseries\small Model sentence frame:}\\[6pt]
    {\small\itshape ``Of the \underline{\hspace{1cm}} [individuals] surveyed, the most
    common [variable] was \underline{\hspace{1.5cm}}, with \underline{\hspace{1cm}}\%
    of responses. The least common was \underline{\hspace{1.5cm}} with
    \underline{\hspace{1cm}}\%.''}

    \vspace{0.45cm}
    {\color{slate}\footnotesize
      A bare ``30\%'' with no context earns no credit. Always anchor your numbers.}
  \end{minipage}
\end{frame}
}

% ─────────────────────────────────────────────────────────────────────────────
% SLIDE 10 — GROUP ACTIVITY
% ─────────────────────────────────────────────────────────────────────────────
{
\setbeamercolor{background canvas}{bg=redbg}
\setbeamercolor{enumerate item}{fg=redacc}
\begin{frame}[t, plain]
  \slideheader{redacc}{Group Activity --- Favorite Sport Survey}
  \begin{minipage}[t]{\textwidth}
    \small

    \textbf{Part 1 (8 min):} Raw data for 30 middle school students.
    \begin{enumerate}
      \setlength{\itemsep}{0.3em}\setlength{\topsep}{0.25em}
      \item Tally the data and build a complete frequency and relative frequency table.
      \item Verify: total $= 30$, sum of relative frequencies $= 1$.
      \item Write 2--3 sentences describing the distribution.
    \end{enumerate}

    \vspace{0.55cm}
    \textbf{Part 2 (5 min):} A second survey of 50 high school students.
    \begin{enumerate}
      \setlength{\itemsep}{0.3em}\setlength{\topsep}{0.25em}
      \item Complete the relative frequency column for the high school data.
      \item Compare the two groups. Which sport is more popular proportionally?
      \item Can you compare using counts alone? Why or why not?
    \end{enumerate}
  \end{minipage}
\end{frame}
}

% ─────────────────────────────────────────────────────────────────────────────
% SLIDE 11 — EXIT TICKET
% ─────────────────────────────────────────────────────────────────────────────
{
\setbeamercolor{background canvas}{bg=sky}
\begin{frame}[t, plain]
  \slideheader{navy}{Exit Ticket --- 5 Minutes}
  \begin{minipage}[t]{\textwidth}
    \small

    A survey of \textbf{50 students} asks: ``Where do you prefer to study?''

    \vspace{0.35cm}
    \setlength{\tabcolsep}{10pt}
    \renewcommand{\arraystretch}{1.35}
    \hspace*{\tabcolsep}\begin{tabular}{@{} >{\bfseries}p{0.24\textwidth}
                        >{\centering\arraybackslash}p{0.18\textwidth}
                        >{\centering\arraybackslash}p{0.22\textwidth} @{}}
      \rowcolor{navy}
      {\color{white}Location} & {\color{white}Count} & {\color{white}Rel.\ freq.} \\
      Library  & 18 & ? \\
      Home     & 22 & ? \\
      Caf\'{e} & 6  & ? \\
      Other    & 4  & ? \\
      \rowcolor{skymid} Total & 50 & ? \\
    \end{tabular}

    \vspace{0.45cm}
    \begin{enumerate}
      \setlength{\itemsep}{0.35em}\setlength{\topsep}{0em}
      \item Complete the relative frequency column.
      \item Write one sentence describing the most common location \textit{in context}.
      \item What percentage of students study somewhere other than home?
    \end{enumerate}
  \end{minipage}
\end{frame}
}

% ─────────────────────────────────────────────────────────────────────────────
% SLIDE 12 — BIG IDEAS / PREVIEW
% ─────────────────────────────────────────────────────────────────────────────
{
\setbeamercolor{background canvas}{bg=navy}
\setbeamercolor{itemize item}{fg=goldacc}
\begin{frame}[t, plain]
  \slideheader{navylight}{Big Ideas \& Preview}
  \begin{minipage}[t]{\textwidth}
    \small

    {\color{goldacc}\bfseries Today's core takeaway:}\\[6pt]
    {\color{white}
    Use \textbf{relative frequency} when comparing groups of different sizes.
    Counts alone are misleading.}

    \vspace{0.55cm}
    {\color{skymid}\bfseries Connections:}
    \begin{itemize}
      \setlength{\itemsep}{0.35em}\setlength{\topsep}{0.25em}
      \item {\color{white}
            \textbf{Lesson 1.4:} the same distributions shown as bar charts and
            pie charts --- tables vs.\ graphs, trade-offs.}
      \item {\color{white}
            \textbf{Unit 2:} two-way tables extend today's one-way tables to
            two categorical variables.}
      \item {\color{white}
            \textbf{Unit 4:} relative frequency $=$ experimental probability ---
            the arithmetic is identical.}
    \end{itemize}

    \vspace{0.55cm}
    {\color{skymid}\small\itshape
      Complete your homework (Lesson 1.3 HW) before next class.}
  \end{minipage}
\end{frame}
}

\end{document}
